% !TEX root = paper.tex
\section{Method}\label{sec:method}

\subsection{Preliminaries}\label{sub:method:preliminaries}

\textbf{Diffusion model basics.}
Diffusion models formulate a forward noising process
over $T$ timesteps for data $\mathbf{x}_0 \sim p_{\text{data}}$
as
\begin{equation}
    q(\mathbf{x}_t \mid \mathbf{x}_{t-1}) = \mathcal{N}\!
    \left(\mathbf{x}_t; \sqrt{1-\beta_t}\mathbf{x}_{t-1}, \beta_t \mathbf{I}\right),
    \label{eq:method:prelim:forward}
\end{equation}
where $\{\beta_t\}$ is a predefined variance schedule.
This admits a closed-form solution
\begin{equation}
    \mathbf{x}_t =
    \sqrt{\bar{\alpha}_t}\mathbf{x}_0 +
    \sqrt{1-\bar{\alpha}_t}\boldsymbol{\epsilon}_t, \quad
    \boldsymbol{\epsilon}_t \sim \mathcal{N}(\mathbf{0}, \mathbf{I}),
    \label{eq:method:prelim:forward-closed}
\end{equation}
where $\bar{\alpha}_t = \prod_{s=1}^t (1-\beta_s)$.
The learned reverse process is characterized by
\begin{equation}
    p_\theta(\mathbf{x}_{t-1} \mid \mathbf{x}_t, \mathbf{c}) = \mathcal{N}\!
    \left(\mathbf{x}_{t-1};
    \boldsymbol{\mu}_\theta(\mathbf{x}_t, t, \mathbf{c}),
    \boldsymbol{\Sigma}_\theta(\mathbf{x}_t, t)\right),
    \label{eq:method:prelim:reverse}
\end{equation}
where $\mathbf{c}$ denotes the conditioning information.
Sampling begins from $\mathbf{x}_T \sim \mathcal{N}(\mathbf{0}, \mathbf{I})$,
iteratively denoising toward $\mathbf{x}_0$.

\textbf{Dynamic Mode Decomposition.}
Dynamic Mode Decomposition (DMD) models the local temporal evolution
of a sequence of snapshots
$\{\mathbf{y}_i\}_{i=1}^m \subset \mathbb{R}^D$
by a linear operator $\mathbf{A} \in \mathbb{R}^{D \times D}$ such that
\[
    \mathbf{y}_{i+1} \approx \mathbf{A}\mathbf{y}_i,
\]
or equivalently,
\[
    \mathbf{Y} \approx \mathbf{A}\mathbf{X}.
\]
Here, $m$ denotes the number of snapshots in the local history window,
and $D$ denotes the dimension of each snapshot vector.
In \Cref{sub:method:dmd-cache},
each $\mathbf{y}_i$ corresponds to a flattened transformer output
at one denoising step.
Given snapshots $\{\mathbf{y}_1, \ldots, \mathbf{y}_m\}$,
the shifted snapshot matrices are defined as
\begin{equation}
	\begin{aligned}
		\mathbf{X} &= [\mathbf{y}_{1}, \ldots, \mathbf{y}_{m-1}]
		\in \mathbb{R}^{D \times (m-1)}, \\
		\mathbf{Y} &= [\mathbf{y}_{2}, \ldots, \mathbf{y}_{m}]
		\in \mathbb{R}^{D \times (m-1)}.
	\end{aligned}
	\label{eq:method:xy}
\end{equation}
A least-squares approximation of $\mathbf{A}$ is given by
\[
    \mathbf{A} \approx \mathbf{Y}\mathbf{X}^{\dagger},
\]
where $(\cdot)^{\dagger}$ denotes the Moore--Penrose pseudoinverse.

The compact singular value decomposition of $\mathbf{X}$ is
\begin{equation}
	\mathbf{X} = \mathbf{U}_r \boldsymbol{\Sigma}_r \mathbf{V}_r^{\top},
	\label{eq:method:svd}
\end{equation}
where
$\mathbf{U}_r \in \mathbb{R}^{D \times r}$,
$\boldsymbol{\Sigma}_r \in \mathbb{R}^{r \times r}$,
$\mathbf{V}_r \in \mathbb{R}^{(m-1) \times r}$,
and $r = \min(D, m-1)$.
The reduced DMD operator
$\tilde{\mathbf{A}} \in \mathbb{R}^{r \times r}$
is then defined as
\begin{equation}
	\tilde{\mathbf{A}} = \mathbf{U}_r^{\top}\mathbf{Y}\mathbf{V}_r\boldsymbol{\Sigma}_r^{-1}.
	\label{eq:method:atilde}
\end{equation}
The eigendecomposition of $\tilde{\mathbf{A}}$ is
\begin{equation}
	\tilde{\mathbf{A}}\mathbf{W} = \mathbf{W}\boldsymbol{\Lambda},
	\label{eq:method:eig}
\end{equation}
where $\mathbf{W} \in \mathbb{C}^{r \times r}$ contains the eigenvectors
and $\boldsymbol{\Lambda} \in \mathbb{C}^{r \times r}$ is the diagonal matrix
of eigenvalues.
The DMD modes
$\mathbf{\Phi} \in \mathbb{C}^{D \times r}$
are given by
\begin{equation}
	\mathbf{\Phi} = \mathbf{Y}\mathbf{V}_r\boldsymbol{\Sigma}_r^{-1}\mathbf{W}.
	\label{eq:method:modes}
\end{equation}
The modal coefficients associated with the latest snapshot are
\begin{equation}
	\mathbf{b} = \mathbf{\Phi}^{\dagger}\mathbf{y}_m.
	\label{eq:method:coeff}
\end{equation}
The one-step prediction is therefore
\begin{equation}
	\hat{\mathbf{y}}_{m+1} = \operatorname{Re}\!\left(\mathbf{\Phi}\boldsymbol{\Lambda}\mathbf{b}\right),
	\label{eq:method:dmd:predict-discrete}
\end{equation}
where $\operatorname{Re}(\cdot)$ denotes the real part.

% \subsection{\texorpdfstring{$D^3$-Cache}{D3-Cache}}\label{sub:method:dmd-cache}

\subsection{\textbf{DMD for Diffusion Model.}}\label{sub:method:dmd-cache}
$D^3$-Cache treats the sequence of diffusion transformer output tensors $\{\mathbf{y}_k\}$
as a temporal signal
and applies DMD
to learn the local linear dynamics
over short temporal windows.
This enables a data-driven predictor
for subsequent transformer outputs,
allowing us to skip expensive forward passes
when recent preview errors indicate reliable extrapolation.

\textbf{Output Representation.}
Diffusion inference proceeds as a discrete-time denoising trajectory
across timesteps $\{t_k\}_{k=1}^K$.
The timestep interval is defined as
\begin{equation}
\Delta t_k = |t_k - t_{k-1}|,
\label{eq:method:dt}
\end{equation}
which may vary across different inference schedules.

At step $k$,
the wrapped diffusion transformer produces a primary output tensor:
\begin{equation}
\mathbf{y}_k = f_\theta(\mathbf{z}_k, t_k, \mathbf{c}) \in \mathbb{R}^D,
\label{eq:method:prelim:diffusion}
\end{equation}
where $\mathbf{z}_k$ is the current latent state,
$\mathbf{c}$ is the conditioning signal (e.g., text embedding),
$\theta$ are the model parameters,
and $D$ is the flattened output dimension.
In implementation,
\(\mathbf{y}_k\) is the primary tensor extracted from the module return value:
the tensor itself,
the \texttt{sample} field,
or the first element of a tuple/list output.

\textbf{Caching Strategy.}
Unlike layer-wise caching methods that cache intermediate activations
$\{\mathbf{h}_k^{(\ell)}\}_{\ell=1}^{L}$ from individual blocks,
$D^3$-Cache operates in the \emph{transformer-output space},
caching only the primary output tensor $\mathbf{y}_k$ returned by the wrapped module.
This choice offers several advantages:
(i) significantly reduced memory footprint compared to storing layer-wise features,
(ii) direct applicability across architectures with different block structures,
(iii) natural alignment with the diffusion sampling process.

We maintain a sliding window history buffer
$\mathcal{H}_k = \{\mathbf{y}_{k-m+1}, \ldots, \mathbf{y}_k\}$
containing the most recent $m$ outputs.
This history is updated only upon full forward pass execution.
Cached predictions are returned to the sampler
but are \emph{not} added to the history,
and we cap the number of consecutive cached steps
to maintain prediction accuracy.


\textbf{History Window and Timestep Statistics.}
We maintain a sliding history window
containing the most recent $m$ transformer outputs:
\begin{equation}
	\mathcal{H}_k = \{ \mathbf{y}_{k-m+1}, \mathbf{y}_{k-m+2}, \ldots, \mathbf{y}_k \},
	\label{eq:method:history}
\end{equation}
and compute the mean timestep interval within this window as
\begin{equation}
	\bar{\Delta t}
	=
	\frac{1}{|\mathcal{D}_k|} \sum_{\Delta t \in \mathcal{D}_k} \Delta t,
	\label{eq:method:avg-dt}
\end{equation}
where $\mathcal{D}_k = \{\Delta t_{k-m+2}, \ldots, \Delta t_k\}$
collects the timestep intervals within the history window.

\textbf{DMD Applied to Diffusion Outputs.}
We treat the sequence of transformer outputs $\{\mathbf{y}_k\}$
as observations of an underlying dynamical system
indexed by the denoising step $k$.
Using the window $\mathcal{H}_k$,
we construct the DMD snapshot matrices $(\mathbf{X}, \mathbf{Y})$
as defined in \Cref{eq:method:xy}
and obtain the DMD components $(\mathbf{\Phi}, \boldsymbol{\Lambda}, \mathbf{b})$
via the procedure in \Cref{sub:method:preliminaries}.

\textbf{Continuous-Time Extrapolation.}
To handle variable timestep spacing in diffusion schedules,
we convert discrete-time eigenvalues to continuous-time frequencies
via
\begin{equation}
	\boldsymbol{\omega} = \frac{1}{\bar{\Delta t}}
	\log\left(\operatorname{diag}(\boldsymbol{\Lambda})\right),
	\label{eq:method:omega}
\end{equation}
where $\log$ is applied element-wise.
Thus $\boldsymbol{\omega} \in \mathbb{C}^r$
in general.
The predicted output for $t_{k+1}$ is then
\begin{equation}
	\hat{\mathbf{y}}_{k+1}(\Delta t_{k+1})
	=
	\text{Re}\left(
		\mathbf{\Phi}
		\left(
			\exp(\boldsymbol{\omega} \Delta t_{k+1})
			\odot
			\mathbf{b}
		\right)
	\right),
	\label{eq:method:predict}
\end{equation}
where $\exp(\cdot)$ is applied element-wise
and $\odot$ denotes the Hadamard product.
This formulation reduces to the standard discrete-time DMD prediction
(Eq.\ref{eq:method:dmd:predict-discrete})
when timestep spacing is uniform.

The predicted output is returned to the sampler
if the adaptive gating criteria
are satisfied; otherwise,
a full forward pass is executed to update the history window.

\subsection{Adaptive Caching Policy}\label{sub:method:adaptive_caching}

\textbf{Preview Error.}
Besides the history buffer,
the cache stores the most recent preview error
computed from an exact forward pass.
Whenever a valid timestep interval is available,
we can form a \emph{preview} prediction
\(\tilde{\mathbf{y}}_{k+1}\)
from the current history
without consuming the skip budget.
After running the exact transformer
and obtaining \(\mathbf{y}_{k+1}\),
we update the gate signal using the batch-averaged relative \( \ell_2 \) error
\begin{equation}
	e_{k+1}
	=
	\frac{1}{B}
	\sum_{b=1}^{B}
	\frac{
		\left\|
			\mathrm{vec}\!\left(\tilde{\mathbf{y}}_{k+1}^{(b)}\right)
			-
			\mathrm{vec}\!\left(\mathbf{y}_{k+1}^{(b)}\right)
		\right\|_2
	}{
		\left\|
			\mathrm{vec}\!\left(\mathbf{y}_{k+1}^{(b)}\right)
		\right\|_2
	},
	\label{eq:method:gate-error}
\end{equation}
where \(B\) is the batch size,
and \(\mathrm{vec}(\cdot)\) flattens each transformer output tensor
into a vector in \(\mathbb{R}^D\).
This matches the implementation,
which compares the preview prediction
against the actual transformer output directly
instead of using auxiliary energy-ratio
or condition-number diagnostics.

\textbf{Adaptive Gating Policy.}
The caching decision is governed by the following hyperparameters:
$S_{\max}$ is the maximum consecutive cached steps,
which prevents error accumulation and long-range drift; and
\(\delta_g\) is the gate threshold applied to the most recent preview error.
$s_k$ is the current skip counter tracking consecutive cached steps
since the last full computation.

At step \(k+1\),
the cache is invoked only when both conditions below are satisfied:
\begin{equation}
	\label{eq:method:gating}
	\begin{aligned}
		e_k &\le \delta_g, \\
		s_k &< S_{\max}
	\end{aligned}
\end{equation}
Otherwise, or whenever the history/timestep state is insufficient,
a full forward pass
$\mathbf{y}_{k+1} = f_\theta(\mathbf{z}_{k+1}, t_{k+1}, \mathbf{c})$
is executed,
the preview error is refreshed using \Cref{eq:method:gate-error},
and the history buffer is updated accordingly.
